% !TEX root = chapter-standalone.tex

\section{Sum}

\subsection{Sum}
\label{sec:disjoint-union}

\begin{ctdefinition}[Sum of two sets]
    \label{def:disjoint-union-of-sets}
    Given sets~\setA and~\setB, their sum (also called \SY{disjoint union}) is the set
    \begin{equation}
        \setA \setdisunion \setB \definedas ( \makeset{ 1 } \cartprod \setA) \setunion ( \makeset{2} \cartprod \setB).
    \end{equation}
\end{ctdefinition}
In other words, an element of~$\setA \setdisunion \setB$ is either a tuple~$\disunionA{\ela}$ for some~$\ela \setin\setA$ or a tuple $\disunionB{\elb}$ for some~$\elb \setin\setB$.
The sets~$\makeset{ 1 }$ and~$\makeset{2}$ here simply provide labels which ``force'' the sets~$\makeset{ 1 } \cartprod \setA$ and~$ \makeset{2} \cartprod \setB$ to be disjoint (even if~\setA and~\setB have elements in common).

%As a simple example, let~$\setA = \makeset{ \sapple, \sgrapes, \scarrot}$ and~$\setB = \makeset{ \stea, \swater }$.
%Then
%\begin{equation}
%    \setA \setdisunion \setB = \makeset{ \disunionA{\sapple}, \disunionA{\sgrapes}, \disunionA{\scarrot}, \disunionB{\stea},  \disunionB{\swater}}.
%\end{equation}

%Or consider~$\setA \setdisunion \setA$, say.
%By definition, this is~$(\makeset{ 1 } \cartprod \setA) \setunion (\makeset{2} \cartprod \setA)$, so if~$\setA = \makeset{ \sapple, \sgrapes, \scarrot}$, then
%\begin{equation}
%    \setA \setdisunion \setA = \makeset{ \disunionA{\sapple}, \disunionA{\sgrapes}, \disunionA{\scarrot}, \disunionB{\sapple}, \disunionB{\sgrapes}, \disunionB{\scarrot} }.
%\end{equation}

Consider the sets~$\setA=\makeset{\sbretzel, \stea}$ and~$\setB=\makeset{\sfondue, \smilk}$.
Their \SY{disjoint union} can be represented as:
%
\equationsag{30_disjoint_union}{eq:fig_disjoint}
%
We can define the \SY{disjoint union} of a set with itself; this corresponds to having two distinct copies of the set:
%
\equationsag{30_disjoint_union_self}{eq:fix_disjointself}
%
Let us also think about what happens when the empty set is at play.
For example, if~$\setA = \Emptyset$, then~$\Emptyset \setdisunion \setB = \makeset{2} \cartprod \setB$, and similarly~$\setA \setdisunion \Emptyset = \makeset{ \disunionA{\ela} \mid \ela \setin\setA }$.
Also,~$\Emptyset \setdisunion \Emptyset = \Emptyset$.

\begin{remark}
    If~\setA and~\setB are finite, then the size of~$\setA \setdisunion \setB$ is the sum of the sizes of~\setA and~\setB:
    \begin{equation}
        \cardof{ \setA \setdisunion \setB} = \cardof \setA + \cardof \setB.
    \end{equation}
    This is a reason why we think of~$\setA \setdisunion \setB$ as a form of addition of sets.
\end{remark}

\begin{remark}
    Later we will see that the \SY{cartesian product} of sets is a special case of a very general construction in category theory, called the \SY{categorical product}, and that the \SY{disjoint union}, or sum, of sets is a special case of a ``dual'' construction, called \SY{categorical coproduct}.
\end{remark}

Analogous to the n-ary cartesian product of any finite number of sets, we can define an n-ary version of disjoint union.

\begin{ctdefinition}[n-ary disjoint union of sets]
    \label{def:n-ary-disjoint-union-of-sets}
    Let $n \setin \natnumbers$.
    Given sets $\setAn{1}, \setAn{2}, \dots, \setAn{n}$, their disjoint union, or sum, is the set
    \begin{equation}
        \setAn{1} + \setAn{2} + \dots + \setAn{n}  \definedas ( \makeset{ 1 } \cartprod \setAn{1}) \setunion ( \makeset{2} \cartprod \setAn{2}) \setunion \dots \setunion ( \makeset{n} \cartprod \setAn{n})
    \end{equation}
\end{ctdefinition}